% !TEX root = ../main.tex
\documentclass[../main.tex]{subfiles}

\begin{document}

\selectlanguage{vietnamese}

\begin{center}
    \textbf{\LARGE BÀI LÀM TỰ LUẬN CHUYÊN ĐỀ ỨNG DỤNG AI} \\
    \vspace{0.2cm}
    \textbf{\large Chủ đề: Multiple Drone Attack in Urban City} \\
    \textit{(Tối ưu hóa chiến lược phòng thủ chống tấn công đa mục tiêu bằng thiết bị bay không người lái trong môi trường đô thị)}
\end{center}

\vspace{0.5cm}

\section{PHẦN 1: ĐẶT VẤN ĐỀ \& MỤC TIÊU (Chiếm 15\% dung lượng)}

\begin{itemize}
    \item \textbf{Bối cảnh thực tế:} Trong bối cảnh an ninh hiện đại, sự gia tăng của các thiết bị bay không người lái (drone) giá rẻ tạo ra mối đe dọa cực lớn đối với an ninh đô thị \cite{mutzari2025defending}. Các tác nhân độc hại có thể phối hợp điều khiển nhiều drone tấn công đồng thời (\textit{Multi-drone attacks}) nhắm vào các hạ tầng trọng yếu như cơ quan chính phủ, nhà máy điện, hệ thống nước và bệnh viện \cite{mutzari2025defending}. Môi trường đô thị với đặc thù nhiều tòa nhà cao tầng, vật cản và nguy cơ nhiễu tín hiệu là một thách thức không thể vượt qua đối với các hệ thống phòng thủ mặt đất truyền thống \cite{mutzari2025defending}. Do drone có lợi thế giá rẻ, cơ động cao, thành phố cần được bảo vệ bằng một phi đội drone đánh chặn (\textit{Defense drones}) hoạt động trên không \cite{mutzari2025defending}.
    \item \textbf{Mục tiêu cốt lõi:} Đồ án này không tiếp cận theo tư duy thụ động, bám đuổi vu vơ ``thấy đâu bắt đó'' của các hệ thống cũ \cite{mutzari2025defending}. Thay vào đó, mục tiêu là ứng dụng AI chiến lược nhằm \underline{tối thiểu hóa thiệt hại về người và tài sản} (\textit{Minimize Damage}) cho toàn thành phố thông qua việc tối ưu hóa phân bổ nguồn lực \cite{mutzari2025defending}. Đồ án chấp nhận một thực tế khách quan rằng tài nguyên phòng thủ của ta luôn bị giới hạn so với số lượng drone của kẻ địch \cite{mutzari2025defending}.
\end{itemize}

\end{document}
